\par\Needspace{7\baselineskip}\bigskip
\begin{center}
\pdfbookmark[2]{Gauss an Bessel, Göttingen, 18. Mai 1814}{pallas-brief-18140518}
\textsc{Gauss} an \textsc{Bessel}, Göttingen, 18. Mai 1814.
\end{center}
\SourceLeaf{430}\noindent\hspace*{\parindent}. . . . .\quad Die Störungen der Pallas durch Jupiter, so weit ich sie berechnen wollte, habe ich vollendet. Die 9 ersten Oppositionen lassen sich mit diesen Störungen allein schon ziemlich gut vereinigen, aber sonderbar genug noch \textls[180]{ungleich} besser, so dass die Fehler im \textls[180]{Durchschnitt} nur \(4''\text{—}5''\) gross werden, \(\BareRoot{\frac{\text{summa quadr. error.}}{\text{multitudo}}}\mathpunct{\raisebox{.45ex}{,}}\) wenn man die Jupitersmasse nicht unbedeutend vergrössert; besser wird sich hierüber urtheilen lassen, wenn erst die Saturnsstörungen mit zugezogen werden, die auch schon vollendet sind. Ich habe jetzt auch die Marsstörungen angefangen, die freilich zwar zahlreich aber fast alle sehr klein sind. Doch ist eine von \(10''\) dabei, die in Zukunft das beste Mittel zur Bestimmung der Marsmasse geben wird.\quad . . . . .
\par
